\section*{Chapitre \ref{ch:b1:matrices} --- Matrices}

\begin{solution}{exo:b1:matrices:1}
\[
AB = \begin{pmatrix} 2 & 1\\ 1 & 0\end{pmatrix},
\quad
BA = \begin{pmatrix} 0 & 1\\ 1 & 2\end{pmatrix},
\quad
A^2 - B^2 = \begin{pmatrix} 1 & 4\\ 0 & 1\end{pmatrix} - I
= \begin{pmatrix} 0 & 4\\ 0 & 0\end{pmatrix},
\]
\[
(A+B)(A-B) = A^2 - AB + BA - B^2
= \begin{pmatrix} 0 & 4\\ 0 & 0\end{pmatrix} +
\begin{pmatrix} -2 & 0\\ 0 & 2 \end{pmatrix}
= \begin{pmatrix} -2 & 4\\ 0 & 2\end{pmatrix}.
\]
Ils diffèrent de $BA - AB \neq 0$~: l'identité $(a+b)(a-b) = a^2 -
b^2$ exige la commutativité, qui est ici en défaut.
\end{solution}

\begin{solution}{exo:b1:matrices:2}
$M$~: $L_2 \leftarrow L_2 - 2L_1$ annule la deuxième ligne~; $L_3
\leftarrow L_3 - L_1$ donne $(0, -1, -2)$. Deux pivots~:
$\operatorname{rk} M = 2$.

$N$~: $L_3 \leftarrow L_3 - L_1$ donne $(0,1,1,1) = L_2$~; puis $L_3
\leftarrow L_3 - L_2 = 0$. Deux pivots~: $\operatorname{rk} N = 2$.
\end{solution}

\begin{solution}{exo:b1:matrices:3}
Réduisons $(A \mid I_3)$~: $L_2 \leftarrow L_2 - 2L_1$, $L_3
\leftarrow L_3 - L_1$~:
\[
\begin{pmatrix}
1 & 0 & 1 & 1 & 0 & 0\\
0 & 1 & -1 & -2 & 1 & 0\\
0 & 1 & 0 & -1 & 0 & 1
\end{pmatrix}
\xrightarrow{L_3 \leftarrow L_3 - L_2}
\begin{pmatrix}
1 & 0 & 1 & 1 & 0 & 0\\
0 & 1 & -1 & -2 & 1 & 0\\
0 & 0 & 1 & 1 & -1 & 1
\end{pmatrix},
\]
puis $L_1 \leftarrow L_1 - L_3$, $L_2 \leftarrow L_2 + L_3$~:
\[
A^{-1} = \begin{pmatrix}
0 & 1 & -1\\
-1 & 0 & 1\\
1 & -1 & 1
\end{pmatrix}.
\]
\emph{Vérification~:} première ligne de $A$ fois première colonne de
$A^{-1}$~: $1 \cdot 0 + 0\cdot(-1) + 1\cdot 1 = 1$~; fois la deuxième
colonne~: $1 - 0 - 1 = 0$~; fois la troisième~: $-1 + 0 + 1 = 0$.
\end{solution}

\begin{solution}{exo:b1:matrices:4}
$u(1) = 1$, $u(X) = X + 1$, $u(X^2) = X^2 + 2X + 1$~: les colonnes de
\omterm{prop:b1:vspaces:coordinates}{coordonnées} dans $(1, X, X^2)$ donnent
\[
M = \begin{pmatrix}
1 & 1 & 1\\
0 & 1 & 2\\
0 & 0 & 1
\end{pmatrix}.
\]
$u$ est inversible parce qu'il possède l'inverse évident $P \mapsto P(X
- 1)$ (composée de substitutions). Sa matrice s'obtient de la même
façon à partir de $u^{-1}(X^k) = (X-1)^k$~:
\[
M^{-1} = \begin{pmatrix}
1 & -1 & 1\\
0 & 1 & -2\\
0 & 0 & 1
\end{pmatrix}.
\]
\end{solution}

\begin{solution}{exo:b1:matrices:5}
$N = \begin{pmatrix} 0 & 1\\ 0 & 0\end{pmatrix}$, $N^2 = 0$. Comme
$2I$ et $N$ commutent, le développement du binôme se tronque après deux
termes~:
\[
A^k = (2I + N)^k = 2^k I + k\,2^{k-1} N
= \begin{pmatrix} 2^k & k\,2^{k-1}\\ 0 & 2^k\end{pmatrix}.
\]
(Vérification pour $k = 2$~: $A^2 = \begin{pmatrix}4 & 4\\ 0 &
4\end{pmatrix}$, ce que confirme le produit direct.)
\end{solution}

\begin{solution}{exo:b1:matrices:6}
Prenons les \omterm{def:b1:matrices:transpose}{traces}~: $\operatorname{tr}(AB - BA) = \operatorname{tr}(AB) -
\operatorname{tr}(BA) = 0$ (\cref{def:b1:matrices:transpose}),
tandis que $\operatorname{tr}(I_n) = n \neq 0$ dans $\R$ ou $\C$. Aucune
solution. (En dimension infinie, l'identité \emph{est} réalisable ---
la dérivation et la multiplication par $x$ la vérifient ---
précisément parce qu'il n'y existe pas de \omterm{def:b1:matrices:transpose}{trace}.)
\end{solution}

\begin{solution}{exo:b1:matrices:7}
Produit télescopique, toutes les puissances de $A$ commutant~:
\[
(I - A)(I + A + \dots + A^{m-1}) = I - A^m = I ,
\]
et la \cref{prop:b1:matrices:ring} promeut l'inverse d'un seul côté en
inverse tout court. Pour l'\omterm{def:b1:logic:map}{application}~: la matrice donnée est $I + N$
avec
\[
N = \begin{pmatrix} 0 & 2 & 3\\ 0 & 0 & 2\\ 0&0&0 \end{pmatrix},
\quad
N^2 = \begin{pmatrix} 0&0&4\\ 0&0&0\\ 0&0&0\end{pmatrix},
\quad N^3 = 0 ,
\]
donc, en remplaçant $A$ par $-N$ dans la formule~:
\[
(I + N)^{-1} = I - N + N^2 =
\begin{pmatrix}
1 & -2 & 1\\
0 & 1 & -2\\
0 & 0 & 1
\end{pmatrix}.
\]
\end{solution}

\begin{solution}{exo:b1:matrices:8}
$A^2 = A$~: l'endomorphisme $a$ est une \omterm{def:b1:linmaps:projection}{projection}
(\cref{thm:b1:linmaps:projchar}), $E = \operatorname{im} a \oplus
\ker a$ avec $\dim\operatorname{im} a = r = \operatorname{rk} A$. Dans
une \omterm{def:b1:vspaces:free}{base} adaptée à cette décomposition ($r$ vecteurs de l'\omterm{def:b1:linmaps:kerim}{image}, puis
une \omterm{def:b1:vspaces:free}{base} du noyau), la matrice de $a$ est
$\begin{pmatrix} I_r & 0\\ 0 & 0\end{pmatrix}$, de \omterm{def:b1:matrices:transpose}{trace} $r$. La
\omterm{def:b1:matrices:transpose}{trace} est invariante par \omterm{def:b1:matrices:changeofbasis}{changement de base}~:
$\operatorname{tr}(P^{-1}MP) = \operatorname{tr}(MPP^{-1}) =
\operatorname{tr} M$ par l'identité cyclique. D'où
$\operatorname{tr} A = r = \operatorname{rk} A$.
\end{solution}

\begin{solution}{exo:b1:matrices:9}
$J^2 = nJ$ (chaque coefficient de $J^2$ somme $n$ fois le nombre $1$).
Cherchons $M^{-1} = \alpha I + \beta J$~:
\[
(aI + bJ)(\alpha I + \beta J)
= a\alpha\, I + (a\beta + b\alpha + nb\beta)\, J .
\]
Ceci vaut $I$ si et seulement si $a\alpha = 1$ et $a\beta + b\alpha +
nb\beta = 0$, c'est-à-dire $\alpha = \frac1a$ et $\beta(a + nb) =
-\frac ba$. Si $a \neq 0$ et $a + nb \neq 0$~:
\[
M^{-1} = \frac 1a I - \frac{b}{a(a + nb)}\, J .
\]
Réciproquement, si $a = 0$~: $M = bJ$ est de rang $\leq 1 < n$ (pour $n
\geq 2$)~: non inversible ($n = 1$ est le cas scalaire). Si $a + nb =
0$~: le vecteur $v = (1, \dots, 1)^{\mathsf T}$ vérifie $Mv = (a +
nb)v = 0$ avec $v \neq 0$~: non inversible. Donc $M \in GL_n \iff a
\neq 0$ et $a + nb \neq 0$.
\end{solution}

\begin{solution}{exo:b1:matrices:10}
\emph{Somme~:} $\operatorname{im}(A + B) \subseteq \operatorname{im} A
+ \operatorname{im} B$ (car $(A+B)x = Ax + Bx$), et la \omterm{thm:b1:findim:grassmann}{formule de
Grassmann} majore la dimension d'une somme par la somme des dimensions.

\emph{Sylvester~:} soit $a$ l'\omterm{def:b1:logic:map}{application} de $A$ restreinte à $V =
\operatorname{im} B$ (de dimension $\operatorname{rk} B$). Son \omterm{def:b1:linmaps:kerim}{image} est
$\operatorname{im}(AB)$ (car $a(Bx) = ABx$), et le \omterm{thm:b1:linmaps:ranknullity}{théorème du rang}
dans $V$ donne~:
\[
\operatorname{rk} B = \dim\ker(a_{|V}) + \operatorname{rk}(AB) .
\]
Or $\ker(a_{|V}) \subseteq \ker A$, de dimension $n -
\operatorname{rk} A$~: donc
\[
\operatorname{rk}(AB) \geq \operatorname{rk} B - (n -
\operatorname{rk} A) = \operatorname{rk} A + \operatorname{rk} B -
n . \qedhere
\]
\end{solution}

\begin{solution}{exo:b1:matrices:11}
\begin{enumerate}
  \item Coefficient par coefficient, $(AD)_{ij} = a_{ij}\,d_j$ et
        $(DA)_{ij} = d_i\,a_{ij}$. Donc $AD = DA$ si et seulement si
        $a_{ij}(d_j - d_i) = 0$ pour tous $i, j$~; lorsque $i \neq j$,
        le facteur $d_j - d_i$ est non nul, ce qui force $a_{ij} = 0$~:
        $A$ est diagonale. Réciproquement, les matrices diagonales
        commutent entre elles.
  \item Si $A$ commute avec toute matrice, elle commute en
        particulier avec $\operatorname{diag}(1, 2, \dots, n)$, donc
        $A = \operatorname{diag}(\lambda_1, \dots, \lambda_n)$ par (1).
        Alors $A E_{ij} = \lambda_i E_{ij}$ (seule la ligne $i$ de
        $E_{ij}$ survit) tandis que $E_{ij} A = \lambda_j E_{ij}$~:
        commuter avec $E_{ij}$ force $\lambda_i = \lambda_j$.
        D'où $A = \lambda I_n$~; et les matrices scalaires commutent
        bien avec tout. Le centre de $\mathcal{M}_n(K)$ est
        $K\,I_n$.
\end{enumerate}
\end{solution}

\begin{solution}{exo:b1:matrices:12}
\begin{enumerate}
  \item Si $\operatorname{rk} A = 1$~: l'\omterm{def:b1:linmaps:kerim}{image} de $A$ est une droite
        $\operatorname{Vect}(C)$, $C \neq 0$, donc la $j$-ième
        colonne de $A$ vaut $\ell_j\,C$ pour des scalaires $\ell_j$
        (non tous nuls), c'est-à-dire $A = C L$ avec $L = (\ell_1,
        \dots, \ell_n) \neq 0$. Réciproquement, si $A = CL \neq 0$,
        toutes les colonnes sont multiples de $C$~: rang $1$.
  \item $A^2 = C\,(L C)\,L$, et $LC$ est le scalaire $\sum_i
        \ell_i c_i = \operatorname{tr}(CL) = \operatorname{tr}
        A$. Donc $A^2 = (\operatorname{tr} A)\,A$, d'où par
        récurrence $A^m = (\operatorname{tr} A)^{m-1} A$. Si
        $\operatorname{tr} A \neq 0$, aucune puissance ne s'annule~;
        si $\operatorname{tr} A = 0$, alors $A^2 = 0$~: une matrice
        de rang $1$ est nilpotente si et seulement si sa \omterm{def:b1:matrices:transpose}{trace} est
        nulle.
  \item Avec $t = \operatorname{tr} A \neq -1$~:
        \[
        (I_n + A)\Bigl(I_n - \frac{A}{1 + t}\Bigr)
        = I_n + A - \frac{A + A^2}{1 + t}
        = I_n + A - \frac{(1 + t)A}{1 + t} = I_n ,
        \]
        en utilisant $A^2 = tA$. Si $t = -1$~: $(I_n + A)A = A + A^2 =
        A - A = 0$ avec $A \neq 0$, donc $I_n + A$ annule toute
        colonne (non nulle) de $A$~: elle n'est pas \omterm{def:b1:logic:inj}{injective}, donc
        pas inversible.
\end{enumerate}
\end{solution}

\begin{solution}{pb:b1:matrices:1}
\textbf{1.} Division euclidienne de $X^n$ par le \omterm{def:b1:poly:def}{polynôme unitaire} $D$
de degré $2$ (\cref{thm:b1:poly:division})~: le quotient et le reste
existent et sont uniques, et le reste est de degré $\leq 1$~: $X^n =
Q_n D + a_n X + b_n$. Pour $n = 0$~: $Q_0 = 0$, $(a_0, b_0) = (0,
1)$~; pour $n = 1$~: $(a_1, b_1) = (1, 0)$.

\textbf{2.} Multiplions par $X$ et réduisons $X^2 = D + sX - p$~:
\[
X^{n+1} = X Q_n D + a_n X^2 + b_n X
= (X Q_n + a_n)\,D + (s\,a_n + b_n)\,X - p\,a_n .
\]
La dernière expression a la forme d'un reste (degré $\leq 1$), donc par
unicité $a_{n+1} = s a_n + b_n$ et $b_{n+1} = -p a_n$.
En substituant $b_{n+1} = -pa_n$ dans $a_{n+2} = s a_{n+1} +
b_{n+1}$, on obtient $a_{n+2} = s\,a_{n+1} - p\,a_n$.

\textbf{3.} Évaluons $X^n = Q_n D + a_n X + b_n$ aux racines~:
$\lambda^n = a_n\lambda + b_n$ et $\mu^n = a_n\mu + b_n$.
En soustrayant et en divisant par $\lambda - \mu \neq 0$~:
\[
a_n = \frac{\lambda^n - \mu^n}{\lambda - \mu},
\qquad
b_n = \lambda^n - a_n\lambda
= \frac{\lambda\mu^n - \mu\lambda^n}{\lambda - \mu} .
\]

\textbf{4.} À la racine double~: $\lambda^n = a_n\lambda + b_n$.
En dérivant l'identité, $nX^{n-1} = Q_n'\,(X - \lambda)^2
+ 2Q_n\,(X - \lambda) + a_n$, et en évaluant en $\lambda$~: $a_n
= n\lambda^{n-1}$~; puis $b_n = \lambda^n - n\lambda^{n} = (1 -
n)\lambda^{n}$.

\textbf{5.} Pour $P = \sum_i p_i X^i$ et $Q = \sum_j q_j X^j$,
\[
P(M)\,Q(M) = \sum_{i,j} p_i q_j M^{i+j} = (PQ)(M),
\]
car les puissances de la seule matrice $M$ commutent entre elles (pour
les sommes, c'est clair par linéarité). Si $D(M) = 0$, la substitution
de $M$ dans $X^n = Q_n D + a_n X + b_n$ donne $M^n = Q_n(M)\,D(M) + a_n
M + b_n I = a_n M + b_n I$.

\textbf{6.} Produits directs~:
\[
A^2 = \begin{pmatrix}
a^2 + bc & b(a + d)\\
c(a + d) & d^2 + bc
\end{pmatrix},
\qquad
s A = \begin{pmatrix}
a(a+d) & b(a+d)\\
c(a+d) & d(a+d)
\end{pmatrix},
\]
donc $A^2 - sA$ a des coefficients hors diagonale nuls et des
coefficients diagonaux valant
$a^2 + bc - a^2 - ad = bc - ad = -p$~: $A^2 - sA + pI_2 = 0$.

\textbf{7.} Avec $A' = \begin{pmatrix} a' & b'\\ c' &
d'\end{pmatrix}$, développons $p(AA') = (aa' + bc')(cb' + dd') -
(ab' + bd')(ca' + dc')$~: les termes $aa'cb'$ et $ab'ca'$ se
simplifient, les termes $bc'dd'$ et $bd'dc'$ se simplifient, et il
reste
\[
aa'dd' - bca'd' + bcb'c' - adb'c'
= (ad - bc)(a'd' - b'c') = p(A)\,p(A').
\]
Si $p \neq 0$, Cayley--Hamilton donne $A\,\bigl(\tfrac1p(sI_2 -
A)\bigr) = \tfrac1p(sA - A^2) = I_2$, d'où l'inverse (et la
\cref{prop:b1:matrices:ring} en fait un inverse des deux côtés). Si $p
= 0$ et si $A$ était inversible, la multiplicativité donnerait $1 =
p(I_2) = p(A)\,p(A^{-1}) = 0$~: impossible. Donc $A \in GL_2 \iff p
\neq 0$.

\textbf{8.} $s = 3$, $p = 2$, $D = X^2 - 3X + 2 = (X - 1)(X -
2)$~: $\lambda = 2$, $\mu = 1$, donc $a_n = 2^n - 1$ et $b_n = 2 -
2^n$ (question 3). D'où
\[
A^n = (2^n - 1)A + (2 - 2^n)I
= \begin{pmatrix} 1 & 2^n - 1\\ 0 & 2^n \end{pmatrix}.
\]
Vérification~: $A^2 = \begin{pmatrix} 1 & 3\\ 0 & 4\end{pmatrix}$, aussi
bien par la formule qu'en élevant directement au carré.

\textbf{9.} $s = 4$, $p = 3\cdot1 - 1\cdot(-1) = 4$~: $D = X^2 -
4X + 4 = (X - 2)^2$, racine double $\lambda = 2$. Question 4~: $a_n
= n\,2^{n-1}$, $b_n = (1 - n)2^n$, donc
\[
A^n = n\,2^{n-1}A + (1 - n)2^n I
= 2^{n-1}\begin{pmatrix} n + 2 & n\\ -n & 2 - n
\end{pmatrix}.
\]
Pour $n = 2$~: $2\begin{pmatrix} 4 & 2\\ -2 & 0\end{pmatrix} =
\begin{pmatrix} 8 & 4\\ -4 & 0 \end{pmatrix}$, qui est bien $A^2$
calculé directement.

\textbf{10.} Par récurrence~: $F^1 = \begin{pmatrix} F_2 & F_1\\ F_1 &
F_0\end{pmatrix}$, et
\[
F^{n+1} = F^n F =
\begin{pmatrix} F_{n+1} + F_n & F_{n+1}\\
F_n + F_{n-1} & F_n \end{pmatrix}
= \begin{pmatrix} F_{n+2} & F_{n+1}\\ F_{n+1} & F_n
\end{pmatrix}.
\]
Ici $s = 1$, $p = -1$, $D = X^2 - X - 1$ de racines $\varphi,
\psi$ ($\varphi - \psi = \sqrt5$, $\varphi\psi = -1$). La
suite $(F_n)$ vérifie $F_0 = 0 = a_0$, $F_1 = 1 = a_1$ et obéit
à la même récurrence que $(a_n)$~: $F_n = a_n = (\varphi^n -
\psi^n)/\sqrt5$, la formule de Binet. Cassini~: en appliquant la
multiplicativité de la question 7 à $F^n$,
\[
F_{n+1}F_{n-1} - F_n^2 = p(F^n) = p(F)^n = (-1)^n .
\]

\textbf{11.} La condition est \omterm{def:b1:linmaps:def}{linéaire} et contient la suite nulle~:
c'est un \omterm{def:b1:vspaces:subspace}{sous-espace vectoriel}. Par récurrence, $u_0, u_1$ déterminent
$u$ linéairement, et tout couple de valeurs initiales est réalisé par
exactement une solution~: comme à l'\cref{exo:b1:findim:10}, $E_D$ est
paramétré bijectivement et linéairement par $(u_0, u_1) \in K^2$~:
$\dim E_D = 2$.

\textbf{12.} $(a_n)$ obéit à la récurrence (question 2) avec $a_0
= 0$, $a_1 = 1$. Il en va de même de $(b_n)$~: $b_{n+2} = -p\,a_{n+1} =
-p(s a_n + b_n) = s\,b_{n+1} - p\,b_n$ (en utilisant deux fois
$b_{n+1} = -pa_n$), avec $b_0 = 1$, $b_1 = 0$. La combinaison $v_n =
u_1 a_n + u_0 b_n$ est alors une solution de valeurs $v_0 = u_0$, $v_1 =
u_1$~; deux solutions de mêmes valeurs initiales coïncident
(récurrence), donc $u_n = u_1 a_n + u_0 b_n$ pour tout $n$.

\textbf{13.} $(\lambda^n)$ est une solution si et seulement si
$\lambda^{n+2} = s\lambda^{n+1} - p\lambda^n$ pour tout $n$,
c'est-à-dire $D(\lambda) = 0$
(après division par $\lambda^n \neq 0$~; notons que $\lambda, \mu \neq
0$ puisque $p = \lambda\mu \neq 0$). Liberté de
$\bigl((\lambda^n), (\mu^n)\bigr)$~: une relation en $n = 0, 1$
donne $c + c' = 0$, $c\lambda + c'\mu = 0$, donc $c(\lambda - \mu)
= 0$~: $c = c' = 0$. Deux vecteurs \omterm{def:b1:vspaces:free}{libres} en dimension $2$~: une \omterm{def:b1:vspaces:free}{base}.
Racine double~: $\bigl((n\lambda^n)\bigr)$ est une solution car,
avec $s = 2\lambda$, $p = \lambda^2$~:
\[
s(n+1)\lambda^{n+1} - p\,n\lambda^n
= \lambda^{n+2}\bigl(2(n+1) - n\bigr) = (n+2)\lambda^{n+2} ;
\]
liberté en $n = 0, 1$~: $c = 0$, puis $c'\lambda = 0$ avec
$\lambda \neq 0$.

\textbf{14.} $D = X^2 - X - 6 = (X - 3)(X + 2)$. Solution générale
$u_n = A\,3^n + B(-2)^n$~; les conditions initiales donnent
$A + B = 1$ et $3A - 2B = 8$, donc $A = 2$, $B = -1$~:
\[
u_n = 2\cdot 3^n - (-2)^n .
\]
Vérification~: $u_2 = 18 - 4 = 14 = u_1 + 6u_0$~; $u_3 = 54 + 8 = 62 =
u_2 + 6u_1 = 14 + 48$.

\textbf{15.} $C\begin{pmatrix} u_n\\ u_{n+1}\end{pmatrix} =
\begin{pmatrix} u_{n+1}\\ -p\,u_n + s\,u_{n+1}\end{pmatrix} =
\begin{pmatrix} u_{n+1}\\ u_{n+2}\end{pmatrix}$, et une récurrence
donne la formule avec $C^n$. De plus $\operatorname{tr} C = 0
+ s = s$ et $p(C) = 0\cdot s - 1\cdot(-p) = p$~: la matrice
compagnon a exactement $D$ pour \omterm{def:b1:poly:def}{polynôme} de Cayley--Hamilton.

\textbf{16.} Pour toute solution $u$ de
\[
u_{n+3} = \alpha u_{n+2} + \beta u_{n+1} + \gamma u_n ,
\]
les vecteurs d'état $v_n = (u_n,
u_{n+1}, u_{n+2})^{\mathsf T}$ vérifient $C_3 v_n = v_{n+1}$
(les deux premières lignes décalent, la dernière applique la
récurrence). Donc
\[
D_3(C_3)\,v_0 = v_3 - \alpha v_2 - \beta v_1 - \gamma v_0 ,
\]
dont les trois composantes valent $u_{k+3} - \alpha u_{k+2} - \beta
u_{k+1} - \gamma u_k = 0$ ($k = 0, 1, 2$). Comme l'état initial
$v_0 = (u_0, u_1, u_2)^{\mathsf T}$ parcourt \emph{tout}
$K^3$ (les valeurs initiales sont \omterm{def:b1:vspaces:free}{libres}), la matrice $D_3(C_3)$ annule
tout vecteur~: $D_3(C_3) = 0$.

\textbf{17.} Écrivons $X^n = Q\,D_3 + R_n$ avec $\deg R_n \leq 2$
et évaluons en chaque racine~: $\lambda_i^n = R_n(\lambda_i)$. Ainsi
$R_n$ est un \omterm{def:b1:poly:def}{polynôme} de degré $\leq 2$ qui interpole les trois
valeurs $\lambda_i^n$ aux trois nœuds distincts $\lambda_i$~: par
unicité dans le \cref{thm:b1:poly:lagrange}, $R_n = \sum_i
\lambda_i^n L_i$ où $(L_i)$ est la \omterm{def:b1:vspaces:free}{base} de Lagrange des nœuds.
En substituant $C_3$ (questions 5 et 16)~:
\[
C_3^{\,n} = R_n(C_3) = \sum_{i=1}^{3} \lambda_i^n\,L_i(C_3),
\]
avec les trois matrices $L_i(C_3)$ indépendantes de $n$~: chaque
coefficient de $C_3^{\,n}$ est une combinaison fixe de $\lambda_1^n,
\lambda_2^n, \lambda_3^n$.

\textbf{18.} $D_3 = X^3 - 2X^2 - X + 2 = (X-1)(X+1)(X-2)$.
Solution générale $u_n = A + B(-1)^n + C\,2^n$. Conditions
initiales~: $A + B + C = 0$, $A - B + 2C = 1$, $A + B + 4C = 1$.
En soustrayant la première de la troisième~: $3C = 1$, $C = \frac13$~;
puis $A + B = -\frac13$ et $A - B = \frac13$~: $A = 0$, $B =
-\frac13$. D'où
\[
u_n = \frac{2^n - (-1)^n}{3}
\]
(les nombres de Jacobsthal). Vérification~: $u_3 = \frac{8 + 1}{3} = 3 =
2u_2 + u_1 - 2u_0 = 2 + 1 - 0$.

\textbf{19.} \omterm{thm:b1:taylor:integral}{Formule de Taylor} du \omterm{def:b1:poly:def}{polynôme} $X^n$ en
$\lambda$~:
\[
X^n = \sum_{k=0}^{n} \binom nk \lambda^{n-k}(X - \lambda)^k ,
\]
et tous les termes avec $k \geq 3$ sont divisibles par $(X -
\lambda)^3$~: le reste est
\[
R_n = \lambda^n + n\lambda^{n-1}(X - \lambda) + \binom
n2\lambda^{n-2}(X - \lambda)^2 .
\]
Pour $M = \lambda I + N$ avec $N^3 = 0$~: $(M - \lambda I)^3 = N^3
= 0$, donc la question 5 donne
\[
M^n = \lambda^n I + n\lambda^{n-1} N + \binom n2
\lambda^{n-2} N^2 ,
\]
ce qui est exactement le développement du binôme de $(\lambda I + N)^n$
tronqué à $N^2$ --- les deux méthodes concordent.

\textbf{20.} L'espace des solutions est de dimension $3$ (même
paramétrage par $(u_0, u_1, u_2)$ qu'à la question 11), et
chaque $(\lambda_i^n)$ est une solution. Liberté~: supposons
$c_1\lambda_1^n + c_2\lambda_2^n + c_3\lambda_3^n = 0$ pour $n =
0, 1, 2$. Fixons $i$ et soit $L_i = \sum_{k \leq 2} p_k X^k$ le
\omterm{def:b1:poly:def}{polynôme} de Lagrange des nœuds vérifiant $L_i(\lambda_j) =
\delta_{ij}$. Alors
\[
0 = \sum_{k=0}^{2} p_k\Bigl(\sum_j c_j\lambda_j^k\Bigr)
= \sum_j c_j\,L_i(\lambda_j) = c_i .
\]
Donc tous les $c_i$ sont nuls~: trois solutions \omterm{def:b1:vspaces:free}{libres} en dimension
$3$, donc une \omterm{def:b1:vspaces:free}{base}~; la solution générale est $c_1\lambda_1^n +
c_2\lambda_2^n + c_3\lambda_3^n$.

\textbf{21.} De $F_k = F_{k+2} - F_{k+1}$, la somme se télescope~:
\[
\sum_{k=1}^{n} F_k = \sum_{k=1}^{n}\bigl(F_{k+2} - F_{k+1}\bigr)
= F_{n+2} - F_2 = F_{n+2} - 1 .
\]

\textbf{22.} Prenons le coefficient $(1,2)$ de $F^{m+n} = F^m F^n$~: le
membre de gauche vaut $F_{m+n}$~; le membre de droite est (la ligne $1$
de $F^m$) fois (la colonne $2$ de $F^n$), c'est-à-dire $F_{m+1}F_n +
F_m F_{n-1}$. Avec $m = n$~:
\[
F_{2n} = F_{n+1}F_n + F_nF_{n-1} = F_n\,(F_{n+1} + F_{n-1}).
\]

\textbf{23.} Par Binet, $F_n - \dfrac{\varphi^n}{\sqrt5} =
-\dfrac{\psi^n}{\sqrt5}$, et $\abs\psi = \frac{\sqrt5 - 1}2 <
1$, donc
\[
\Bigl|F_n - \frac{\varphi^n}{\sqrt5}\Bigr|
\leq \frac{1}{\sqrt5} < \frac12
\qquad (n \geq 0):
\]
$F_n$ est l'entier le plus proche de $\varphi^n/\sqrt5$.

\textbf{24.} $t_n = F_{n+1} + F_{n-1}$ est une combinaison de
suites de Fibonacci décalées, donc vérifie la même
récurrence~: $t_{n+2} = t_{n+1} + t_n$~; et $t_1 = F_2 + F_0 =
1$, $t_2 = F_3 + F_1 = 3$~: ce sont les nombres de Lucas $L_n$.
La suite $\varphi^n + \psi^n$ est une solution de mêmes
deux premières valeurs ($\varphi + \psi = 1$, $\varphi^2 + \psi^2 = (
\varphi + \psi)^2 - 2\varphi\psi = 3$), donc $L_n = \varphi^n +
\psi^n$. Enfin
\[
F_n L_n = \frac{(\varphi^n - \psi^n)(\varphi^n +
\psi^n)}{\sqrt5} = \frac{\varphi^{2n} - \psi^{2n}}{\sqrt5} =
F_{2n},
\]
ce qui retrouve la question 22.

\textbf{25.} (i) Les cinq matrices $I, A, A^2, A^3, A^4$ vivent dans
$\mathcal{M}_2(K)$, de dimension $4$, donc \emph{un certain} \omterm{def:b1:poly:def}{polynôme}
non nul de degré $\leq 4$ annule $A$~; la partie II a précisé cela
en la quadratique explicite $A^2 = sA - pI$, qui enferme toutes les
puissances dans le plan $\operatorname{Vect}(I, A)$. (ii)
La division euclidienne réduit $X^n$ modulo cette quadratique, et les
deux coefficients du reste obéissent à la récurrence à deux termes
$a_{n+2} = s\,a_{n+1} - p\,a_n$~: l'exponentiation est devenue une
itération. (iii) La question 6 est le cas $n = 2$ du
\emph{théorème de Cayley--Hamilton}, valable en toute dimension et
démontré dans le volume de Licence~2. (iv) La matrice compagnon boucle
la boucle~: toute \omterm{pb:b1:matrices:1}{récurrence linéaire} \emph{est} une puissance de
matrice, le même \omterm{def:b1:poly:def}{polynôme} $D$ apparaissant comme données de \omterm{def:b1:matrices:transpose}{trace} et de
déterminant, de sorte que le calcul des restes résout les récurrences
et calcule les puissances d'un seul geste.
\end{solution}
